%% kapa: exam=1 round=19 year=2008 period=1 form=A subject=accounting subject_ko=회계학 seq=28 number=108
%% source: Q-Net 감정평가사 기출문제 (cst00309) artlSeq=1001962 file=extracted/1008/1001962/2008년 제19회 감정평가사 1차시험 기출문제/1교시.hwp
%% conversion: kordoc 4.24.51 -> tools/gen1.py (auto); review=proofread-sonnet; flags=none
\begin{question}{108}
(주)설악전자는 동부와 중부, 서부 세 개의 지역사업부를 가지고 있다. 20X1년 각 사업부의 수익과 비용관련 정보는 다음과 같다.

{\small\centering
\begin{tabular}{lrrr}
\hline
 & \multicolumn{1}{c}{\uline{동\hspace{1em}부}} & \multicolumn{1}{c}{\uline{중\hspace{1em}부}} & \multicolumn{1}{c}{\uline{서\hspace{1em}부}} \\
매출 & ￦720,000 & ￦570,000 & ￦350,000 \\
변동원가 & \uline{370,000} & \uline{290,000} & \uline{240,000} \\
공헌이익 & 350,000 & 280,000 & 110,000 \\
추적고정원가 & 130,000 & 90,000 & 80,000 \\
공통고정원가 & \uline{120,000} & \uline{70,000} & \uline{50,000} \\
영업이익(손실) & \underline{\underline{￦100,000}} & \underline{\underline{￦120,000}} & \underline{\underline{￦(20,000)}} \\
\hline
\end{tabular}
\par}

(주)설악전자의 경영진에서는 손실을 보고한 서부사업부 폐지를 검토하고 있다. 서부사업부가 폐지된다면 서부사업부의 추적고정원가는 발생하지 않을 것으로 조사되었다. 그러나 공통고정원가는 영향을 받지 않는다. 이러한 자료를 토대로 할 때 서부사업부를 폐지한다면 (주)설악전자의 영업이익은 얼마가 될 것인가? \par
\begin{choices}
  \item ￦170,000
  \item ￦190,000
  \item ￦200,000
  \item ￦220,000
  \item ￦240,000
\end{choices}
\end{question}
\answer{1}
